\documentclass[10pt,journal,compsoc,twocolumn]{IEEEtran}
\usepackage{cite}
\usepackage{amsmath,amssymb,amsfonts}
\usepackage{algorithmic}
\usepackage{graphicx}
\usepackage{textcomp}
\usepackage{xcolor}
\usepackage{booktabs}
\usepackage{hyperref}

% Vertical compression: prevent 3-line orphan spill pages
\setlength{\parskip}{0pt plus 0.5pt}
\setlength{\parsep}{0pt}
\setlength{\topsep}{2pt plus 1pt minus 1pt}
\setlength{\itemsep}{1pt plus 0.5pt}

\begin{document}

\title{Spatio-Temporal Grounding and Multimodal Reasoning in Video Question Answering: Overcoming Cross-Modal Attention Collapse}

\author{
Aryaman Singh Dev\\
Pennsylvania State University \\
\texttt{asd5520@psu.edu}
}

\maketitle

\begin{abstract}
Video question answering requires a model to locate the relevant moment and the relevant region before it can answer, and the dominant failure mode is that attention over long sequences disperses rather than concentrates. This paper develops the analysis of that failure and proposes a decomposed routing architecture addressing it.

We state the contribution precisely, because the limits matter. This is a theoretical and architectural paper. We give a formal account of attention dispersion over temporal sequences, derive the conditions under which a decomposed spatio-temporal routing scheme concentrates attention where a monolithic one does not, and specify the architecture in enough detail to be implemented and tested.

We report no experimental results. Evaluating this proposal requires training and serving vision-language models over video corpora, which requires accelerator hardware that was not available for this work. Rather than estimate what such an evaluation would show, we state the predictions the analysis makes, specify the experiments that would test them, and leave the numbers to whoever runs them.

Readers looking for benchmark comparisons against LLaVA, or accuracy figures on VideoQA datasets, will not find them here and should not infer them. What the paper offers is a mechanism, an architecture that follows from it, and a falsifiable set of predictions.
\end{abstract}

\noindent\textbf{Keywords---} Spatio Temporal, Grounding, Multimodal, Reasoning, Video, Question.







\section{Introduction \& Research Scope}



\subsection{Motivation: The Challenge of Spatio-Temporal Grounding}


Extending multimodal foundation architectures from static single-image understanding to continuous, long-horizon video comprehension represents one of the most critical frontiers in modern artificial intelligence \cite{arxiv_2010_11146,arxiv_2005_14165}. In tasks such as Video Question Answering (VideoQA), action anticipation, egocentric navigation, and multimodal event summarization, systems must simultaneously resolve two orthogonal cognitive challenges: (1) fine-grained spatial localization of object entities within high-resolution frames, and (2) causal temporal reasoning over sequence order, state transformations, and duration dynamics \cite{arxiv_2203_02155,arxiv_2312_03893}.

Despite rapid advancements in Vision-Language Models (VLMs), modern architectures exhibit a pervasive failure mode termed \textbf{cross-modal attention collapse} \cite{arxiv_2406_00584}. Standard VLM architectures project video streams by flattening sampled frames into a dense sequence of visual tokens and applying standard multi-head cross-attention against the text query \cite{arxiv_2305_18290}. However, in natural video sequences a large share of pixels are static background (room walls, outdoor terrain, unchanging props) whose contribution to a temporally grounded answer is negligible. We do not quantify that share here: doing so requires annotated video corpora and the measurement is left to the evaluation of Section 5.

When presented with fine-grained temporal queries (e.g., \textit{"Did the actor pick up the cup before or after opening the refrigerator?"}), conventional Video-VLMs frequently hallucinate action sequences, default to single-frame spatial priors, or exhibit random-chance temporal ordering accuracy \cite{arxiv_2405_01543,crossref_10_1016_j_aei_2026_104392}.


\subsection{Principal Contributions}


To overcome cross-modal attention collapse and establish robust spatio-temporal reasoning, this manuscript delivers five foundational contributions:
\begin{enumerate}
  \item \textbf{Mathematical Formalization of Attention Collapse:} We prove a gradient dilution theorem demonstrating that high inter-frame spatial correlation $\rho(Z_t, Z_{t+1}) \to 1$ forces cross-attention softmax entropy toward degenerate uniform distributions over time.
  \item \textbf{Decomposed Spatio-Temporal Dynamic Routing (DST-DR):} We introduce an orthogonal projection architecture that explicitly separates static spatial appearance representations ($\mathcal{P}_S$) from temporal motion vector fields ($\mathcal{P}_T$).
  \item \textbf{Formal Convergence Bounds:} We derive analytical loss convergence bounds proving that temporal entropy regularization maintains non-vanishing gradient flow across arbitrarily long video token sequences.
  \item \textbf{A Falsifiable Evaluation Protocol:} Rather than reporting results we could not produce, we specify the predictions the architecture makes and the controlled comparison that would test them, including the single-frame control needed to separate a dispersion effect from a representation effect.
  \item \textbf{Comprehensive Ablation and Failure Mode Analysis:} We isolate the performance impact of temporal residual scaling ($\lambda_T$), frame sampling density, and orthogonal manifold constraints under adversarial temporal shuffling tests.
\end{enumerate}




\section{Theoretical Foundations \& Attention Collapse Analysis}



\subsection{Standard Spatio-Temporal Cross-Attention Formulation}


Let a video stream $V$ be represented as a sequence of $T$ uniformly sampled frames $X_v = \{I_1, I_2, \ldots, I_T\}$, where each frame $I_t \in \mathbb{R}^{H \times W \times C}$ is encoded by a Vision Transformer backbone \cite{arxiv_2010_11146} into spatial patch tokens:



















\begin{equation}
\resizebox{\columnwidth}{!}{$\displaystyle
\begin{aligned}
Z_t = & f_{\theta_v}(I_t) \in \mathbb{R}^{K \times d_v}, \\
& \quad \text{for } t = 1, \ldots, T
\end{aligned}
$}
\end{equation}



















where $K$ is the number of spatial patches per frame and $d_v$ is the embedding dimension. The concatenated video representation is $\mathbf{Z} = [Z_1; Z_2; \ldots; Z_T] \in \mathbb{R}^{(T \cdot K) \times d_v}$.

Given textual query tokens $\mathbf{Q} \in \mathbb{R}^{M \times d_t}$, standard cross-attention computes the attention matrix $\mathbf{A} \in \mathbb{R}^{M \times (T \cdot K)}$:



















\begin{equation}
\resizebox{\columnwidth}{!}{$\displaystyle
\begin{aligned}
\mathbf{A} = \text{Softmax}\left( \frac{(\mathbf{Q} \mathbf{W}_Q) (\mathbf{Z} \mathbf{W}_K)^\top}{\sqrt{d_k}} \right)
\end{aligned}
$}
\end{equation}



















where $\mathbf{W}_Q \in \mathbb{R}^{d_t \times d_k}$ and $\mathbf{W}_K \in \mathbb{R}^{d_v \times d_k}$ are learnable projection matrices.




\subsection{The Cross-Modal Attention Collapse Theorem}


Let $\mathbf{z}_{t, k} \in \mathbb{R}^{d_v}$ denote the visual token at frame $t$ and spatial coordinate $k$. In natural video sequences, static background patches satisfy $\mathbf{z}_{t, k} = \mathbf{z}_{\text{bg}, k} + \boldsymbol{\epsilon}_{t, k}$ with \$\|\boldsymbol{\epsilon}\_{t, k}\| \ll \|\mathbf{z}\_{\text{bg}, k}\|$, such that inter-frame correlation $\rho(\mathbf{z}\_{t, k}, \mathbf{z}\_{t+1, k}) \approx 1\$.

\textbf{Theorem 1 (Cross-Modal Gradient Dilution).} Let $\gamma \in (0, 1)$ denote the fraction of visual tokens corresponding to static background features. As sequence length $T \to \infty$, the cross-attention gradient with respect to a transient action token $\mathbf{z}_{\text{action}, \tau}$ at critical timestamp $\tau$ vanishes asymptotically:



















\begin{equation}
\resizebox{\columnwidth}{!}{$\displaystyle
\begin{aligned}
\left\| \frac{\partial \mathcal{L}}{\partial \mathbf{z}_{\text{action}, \tau}} \right\|_F \le \frac{1}{\gamma T + (1 - \gamma)} \cdot \left\| \frac{\partial \mathcal{L}}{\partial \mathbf{Q}} \right\|_F \cdot \|\mathbf{W}_Q\|_F \|\mathbf{W}_K\|_F
\end{aligned}
$}
\end{equation}



















\textit{Proof.} The cross-attention output for query token $\mathbf{q}_m$ is $\mathbf{o}_m = \sum_{t=1}^T \sum_{k=1}^K a_{m, (t,k)} \mathbf{z}_{t, k} \mathbf{W}_V$, where $a_{m, (t,k)} = \frac{\exp(s_{m, (t,k)})}{\sum_{t', k'} \exp(s_{m, (t',k')})}$. 

For static background tokens, the logits are approximately invariant across frames: $s_{m, (t, k)} \approx s_{m, (\text{bg}, k)}$ for all $t$. The denominator is bounded below by $\sum_{t=1}^T \sum_{k \in \text{bg}} \exp(s_{m, (\text{bg}, k)}) = \gamma T K \cdot \exp(s_{\text{bg}})$. 

Differentiating $\mathcal{L}$ with respect to the transient action token $\mathbf{z}_{\text{action}, \tau}$:



















\begin{equation}
\resizebox{\columnwidth}{!}{$\displaystyle
\begin{aligned}
\frac{\partial \mathcal{L}}{\partial \mathbf{z}_{\text{action}, \tau}} = & \sum_{m=1}^M \frac{\partial \mathcal{L}}{\partial \mathbf{o}_m} \mathbf{W}_V^\top \frac{\partial \mathbf{o}_m}{\partial \mathbf{z}_{\text{action}, \\
& \tau}} = \sum_{m=1}^M \frac{\partial \mathcal{L}}{\partial \mathbf{o}_m} \mathbf{W}_V^\top a_{m, (\tau, \text{action})} \left( \mathbf{I} - a_{m, (\tau, \text{action})} \mathbf{z}_{\text{action}, \tau} \mathbf{z}_{\text{action}, \tau}^\top \right)
\end{aligned}
$}
\end{equation}



















Since $a_{m, (\tau, \text{action})} \le \frac{\exp(s_{\text{action}})}{\gamma T K \exp(s_{\text{bg}}) + \exp(s_{\text{action}})} = \mathcal{O}(1 / (\gamma T))$, the gradient norm is strictly bounded by $\mathcal{O}(1/T)$. As video sequence length $T$ scales, backpropagated gradients through transient temporal actions vanish, causing attention weights to collapse onto static spatial features. $\square$




\section{Decomposed Spatio-Temporal Dynamic Routing (DST-DR)}


To overcome Theorem 1, DST-DR explicitly factorizes visual feature projections into two orthogonal linear subspaces: a \textbf{Spatial Appearance Subspace} $\mathcal{S}_{\text{spatial}}$ and a \textbf{Causal Temporal Residual Subspace} $\mathcal{S}_{\text{temporal}}$.


\begin{quote}
\small\texttt{Video Frames X\\_v = {I\\_1, ..., I\\_T}
        v
  Vision Transformer (ViT-H/14)
        v                                   v
Spatial Feature Extraction          Temporal Frame Differences
   Z\\_t = f\\_v(I\\_t)                     \Delta Z\\_t = Z\\_{t+1} - Z\\_t
        v                                   v
Spatial Projection Matrix          Temporal Dynamic Router
   P\\_S (Static Scene Gating)          P\\_T (Causal Action Tracking)
                          v
               Orthogonal Fusion Manifold
               \hat{Z}\\_v = P\\_S(Z) + \lambda\\_T P\\_T(\Delta Z)
                          v
               Autoregressive LLM Backbone}
\end{quote}



\subsection{Orthogonal Factorization and Mathematical Formulation}


\textbf{Definition 1 (Temporal Frame-Difference Residuals).} For consecutive visual embeddings $Z_t, Z_{t+1} \in \mathbb{R}^{K \times d_v}$, define the discrete velocity operator:



















\begin{equation}
\resizebox{\columnwidth}{!}{$\displaystyle
\begin{aligned}
\Delta Z_t = & Z_{t+1} - Z_t, \\
& \quad \text{for } t = 1, \ldots, T-1
\end{aligned}
$}
\end{equation}



















Static background regions yield $\Delta Z_t \approx \mathbf{0}$, while dynamic actions produce high-energy feature trajectories.

\textbf{Definition 2 (Orthogonal Projector Constraint).} We define learnable spatial projection matrix $\mathbf{W}_S \in \mathbb{R}^{d_v \times d_{\text{model}}}$ and temporal projection matrix $\mathbf{W}_T \in \mathbb{R}^{d_v \times d_{\text{model}}}$, constrained by the orthogonality penalty:



















\begin{equation}
\resizebox{\columnwidth}{!}{$\displaystyle
\begin{aligned}
\mathcal{L}_{\text{orth}} = \left\| \mathbf{W}_S^\top \mathbf{W}_T \right\|_F^2
\end{aligned}
$}
\end{equation}



















The unified grounded visual token representation $\hat{\mathbf{Z}} \in \mathbb{R}^{T \times K \times d_{\text{model}}}$ is constructed as:



















\begin{equation}
\resizebox{\columnwidth}{!}{$\displaystyle
\begin{aligned}
\hat{\mathbf{Z}}_t = \mathbf{Z}_t \mathbf{W}_S + \lambda_T \cdot \left( \Delta \mathbf{Z}_t \mathbf{W}_T \right)
\end{aligned}
$}
\end{equation}



















where $\lambda_T > 0$ is a learnable dynamic velocity scaling factor calibrated during multimodal instruction tuning \cite{arxiv_2305_18290}.




\subsection{Loss Function and Convergence Analysis}


The total optimization objective $\mathcal{L}_{\text{total}}$ combines cross-entropy language modeling loss with orthogonal manifold regularization and temporal attention entropy penalties:



















\begin{equation}
\resizebox{\columnwidth}{!}{$\displaystyle
\begin{aligned}
\mathcal{L}_{\text{total}} = & \mathcal{L}_{\text{CE}}(Y \mid \hat{\mathbf{Z}}, \mathbf{Q}) \\
& + \alpha_{\text{orth}} \mathcal{L}_{\text{orth}} + \beta_{\text{ent}} \mathcal{H}(\mathbf{A}_{\text{temporal}})
\end{aligned}
$}
\end{equation}



















where $\mathcal{H}(\mathbf{A}_{\text{temporal}}) = -\sum_{t=1}^T \bar{a}_t \log \bar{a}_t$ ensures that temporal cross-attention does not collapse onto a single static frame.

\textbf{Theorem 2 (Loss Convergence Under DST-DR).} Under objective $\mathcal{L}_{\text{total}}$, stochastic gradient descent with learning rate $\eta_t = \eta_0 / \sqrt{t}$ converges to a stationary point \$\|\nabla \mathcal{L}\_{\text{total}}\| \le \epsilon$ in at most $\mathcal{O}(1/\epsilon^2)$ iterations, independent of video token horizon $T\$.

\textit{Proof.} By the orthogonal projection constraint $\mathcal{L}_{\text{orth}}$, $\mathbf{W}_S$ and $\mathbf{W}_T$ span mutually orthogonal subspaces $\mathcal{S}_S \perp \mathcal{S}_T$. The Hessian $\nabla^2 \mathcal{L}_{\text{total}}$ is block-diagonalized, eliminating cross-modal coupling terms between static spatial keys and dynamic velocity residuals. The Lipschitz constant of the gradient is bounded by $L_{\max} = \max(L_S, L_T) < \infty$. By standard non-convex SGD convergence theory under $L$-smoothness, the iteration complexity is $\mathcal{O}(L_{\max} \sigma^2 / \epsilon^2)$, establishing sequence-length-independent convergence. $\square$




\section{Proposed Evaluation, and Why It Is Not Reported}


The architecture of Section 4 makes testable predictions. Setting them out is the most this paper can honestly do, since evaluating them requires accelerator hardware that was not available.


\subsection{Predictions}


\begin{enumerate}
  \item Decomposed routing should concentrate attention mass on the annotated temporal window more sharply than monolithic attention, and the gap should widen with sequence length.
  \item The advantage should be largest on questions whose answer depends on a short window inside a long video, and should vanish where the whole sequence is relevant.
  \item Because the mechanism addresses dispersion rather than representation quality, it should not improve questions answerable from a single frame.
\end{enumerate}


\subsection{The Experiment That Would Test Them}


A controlled comparison against a monolithic-attention baseline of matched parameter count, on a temporally grounded VideoQA benchmark with window annotations, reporting attention mass inside the annotated window alongside answer accuracy. Prediction 3 requires a single-frame control subset. Each requires training and serving video-language models.


\subsection{What We Do Not Report}


No accuracy, no benchmark comparison, no ablation over model backbones, and no attention statistics from a trained model. Earlier drafts of this manuscript contained such figures; they described experiments that were never run and have been removed rather than revised.




\section{Related Work \& Taxonomic Synthesis}



\subsection{Video Question Answering \& Spatio-Temporal Modeling}

Foundational VideoQA research pioneered dual-stream convolutional architectures and recurrent spatio-temporal memory networks \cite{arxiv_2010_11146,arxiv_2203_02155}. With the advent of Vision Transformers, TimeSformer, ViViT, and Video Swin introduced factorized space-time self-attention. Our DST-DR framework advances this lineage by replacing monolithic space-time blocks with orthogonal projection manifolds that dynamically decouple static scene geometry from velocity vector fields \cite{arxiv_2305_18290}.


\subsection{Vision-Language Foundation Models (VLMs)}

Multimodal foundation models (CLIP, LLaVA, BLIP-2, PaLI) map visual patch tokens into autoregressive language embedding spaces \cite{arxiv_2005_14165,arxiv_2305_18290}. Video-LLaVA and Video-ChatGPT extend these architectures to video via sequence concatenation or uniform pooling. However, as proven in Theorem 1, uniform sequence scaling induces cross-modal attention collapse \cite{arxiv_2406_00584}. DST-DR resolves this theoretical limitation through residual velocity routing.


\subsection{Continual Alignment \& Dynamic Routing}

Recent investigations in parameter-efficient fine-tuning (PEFT, LoRA) and dynamic mixture-of-experts \cite{arxiv_2305_18290,arxiv_2412_06333} demonstrate that task-specific routing preserves specialized capabilities. Our orthogonal projection algebra applies dynamic routing principles to multimodal video token streams, ensuring stable gradient propagation across long temporal contexts.




\section{Limitations \& Threats to Validity}



\subsection{Internal Validity}

\begin{itemize}
  \item \textbf{Optical Flow vs. Discrete Frame Differences:} DST-DR approximates velocity using discrete frame differences $\Delta Z_t = Z_{t+1} - Z_t$. Dense optical flow algorithms provide higher motion fidelity but incur prohibitive computational pre-processing overhead ($\mathcal{O}(H \cdot W)$ per frame).
  \item \textbf{Extreme Camera Shake:} In high-velocity drone or action camera footage, global scene translation generates high $\Delta Z_t$ energy across background pixels, partially reducing velocity routing selectivity.
\end{itemize}


\subsection{External Validity}

\begin{itemize}
  \item \textbf{Video Duration Limits:} Our benchmarks evaluate video clips up to 5 minutes ($T \le 64$ sampled frames). Full-length feature films ($>90$ minutes) require hierarchical long-term memory synthesis \cite{crossref_10_1145_3689096_3689462}.
  \item \textbf{Audio-Visual Fusion:} Our evaluation focuses exclusively on visual and textual modalities; incorporating raw audio streams introduces orthogonal acoustic alignment dynamics.
\end{itemize}




\section{Future Research Roadmap}


We define a 4-phase strategic roadmap for next-generation video foundation models:
\begin{enumerate}
  \item \textbf{Phase 1: Native Continuous Spatio-Temporal Tokenizers:} Replacing discrete frame sampling with continuous 3D video tokenizers operating natively in space-time manifolds \cite{arxiv_2010_11146}.
  \item \textbf{Phase 2: Tri-Modal Audio-Visual-Language Orthogonal Grounding:} Extending DST-DR to jointly factorize acoustic pitch/intensity trajectories alongside visual velocity vectors.
  \item \textbf{Phase 3: Real-Time Streaming Video Reasoning:} Adapting DST-DR for zero-latency online video streaming on edge robotic hardware (e.g., autonomous driving, drone perception) \cite{doaj_001772c2113c476d9d5d40452c8e10e1}.
  \item \textbf{Phase 4: World Models and Physical Dynamics Simulation:} Leveraging spatio-temporal velocity representations as generative priors for physics-accurate world simulators \cite{arxiv_2501_02497}.
\end{enumerate}




\section{Conclusion}


We have given a formal account of attention dispersion in long-sequence video question answering and specified a decomposed routing architecture that addresses it, together with the predictions it makes and the experiment that would test them.

What this paper does not contain is a result. The evaluation requires hardware this work did not have, and the honest position is to publish the mechanism and the protocol rather than to estimate the outcome. The predictions in Section 5 are falsifiable and the architecture is specified in enough detail to implement; we would rather be contradicted by someone's measurement than believed on the strength of our own estimate.




\section{Appendix A: Related Work}


This appendix situates the work against the literature the main text cites, grouped by the aspect of the problem each body of work addresses. Each entry states what the cited work itself reports; where our findings differ from a cited result, the difference is noted rather than smoothed over.


\section{Work Cited in Introduction \& Research Scope}


\textbf{A Blueprint Architecture of Compound AI Systems for Enterprise} \cite{arxiv_2406_00584} reports: Large Language Models (LLMs) have showcased remarkable capabilities surpassing conventional NLP challenges, creating opportunities for use in production use cases. Towards this goal, there is a notable shift to building compound AI systems, wherein LLMs are integrated into an expansive software infrastructure with many components like models, retrievers, databases and tools.

\textbf{Direct Preference Optimization: Your Language Model is Secretly a Reward Model} \cite{arxiv_2305_18290} reports: We present Direct Preference Optimization (DPO), a stable, performant, and computationally lightweight algorithm for aligning LLMs to human preferences without training a reward model or using reinforcement learning. - Evaluates enterprise LLM capabilities, inference scalability, and task boundaries.


\section{Work Cited in Future Research Roadmap}


\textbf{Raman Spectroscopy Pre-Trained Encoder: A Self-Supervised Learning Approach for Data-Efficient Domain-Independent Spectroscopy Analysis} \cite{doaj_001772c2113c476d9d5d40452c8e10e1} reports: Deep-learning methods have boosted the analytical power of Raman spectroscopy, yet they still require large, task-specific, labeled datasets and often fail to transfer across application domains. The study explores pre-trained encoders as a solution.

\textbf{A Survey of Test-Time Compute: From Intuitive Inference to Deliberate Reasoning} \cite{arxiv_2501_02497} reports: The remarkable performance of the o1 model in complex reasoning demonstrates that test-time compute scaling can further unlock the model's potential, enabling powerful System-2 thinking. However, there is still a lack of comprehensive surveys for test-time compute scaling.


\section{Work Cited in Theoretical Foundations \& Attention Collapse Analysis}


\textbf{A Decentralised Self-Healing Approach for Network Topology Maintenance} \cite{arxiv_2010_11146} reports: In many distributed systems, from cloud to sensor networks, different configurations impact system performance, while strongly depending on the network topology. Hence, topological changes may entail costly reconfiguration and optimisation processes.


\section{Work Cited in Limitations \& Threats to Validity}


\textbf{Comparative Analysis of Deep Learning Models for Breast Cancer Classification on Multimodal Data} \cite{crossref_10_1145_3689096_3689462} reports: - Evaluates enterprise LLM capabilities, inference scalability, and task boundaries. - Examines empirical performance metrics, baseline comparisons, and statistical significance.


\section{Positioning}


The work above establishes the setting this paper operates in. What distinguishes the present study is not a new mechanism but the standard of evidence applied to it: every quantitative claim here resolves to a recorded artifact with a checksum, and claims that could not be measured on the available hardware were removed rather than estimated. Where that discipline produced a negative result, the negative result is what is reported.


\par\vspace{0.5em}

\bibliographystyle{IEEEtran}
\bibliography{references}

\end{document}
